% ----------------------------------------------



\chapter{Introduction}\label{ch:structure}

Ce projet porte sur la modélisation de la distance de décollage d’un avion à l’aide d’un modèle simplifié basé sur les équations fondamentales de la mécanique du vol.

L’objectif est d’identifier les principaux paramètres influençant la distance de décollage et de proposer un modèle de simulation permettant d’estimer cette distance dans différentes conditions.

\vspace{1cm}
%============================================================================================================
%
%================================================================================================================

\section{Starter to be modify} 

	\subsection{Compétence à acquérir}
	L'analyse de ce rapport permettra de développer les compétences suivantes :
	\begin{itemize}[itemsep=0pt]
		\item \textbf{Comprendre :} la disposition des force (comment elle affecte notre avion ) lors d'une montée stable et importante de l'angle de montée ($\gamma$)
		\item \textbf{Calculer  :} la vitesse ascensionnelle(Vc) en utilisant la notion de puissance excédentaire
		\item \textbf{l'influence :} de la densité de l'air et l'altitude sur les performances des turboréacteur (aussi des avion a hélice )
		\item \textbf{Déterminer  :} graphiquement et analytiquement les plafond de service et d'un avion
	\end{itemize}

	
	\subsection{Outils utiliser pour ce rapport}
	
		Pour mener cette étude ,nous avons utiliser 
		\begin{itemize}
			\item  Les modèles mathématiques de performance de montée proviennent du cours de mécanique du vol, issu de différent livre.
			\item Nous utilisons des applications comme VS Code pour la programmation python des codes afin de réaliser des analyses graphiques (courbes de poussée et de puissance) et des calculs numériques pour résoudre des équations  telle que celle de ratio de vitesse (Rv).
		\end{itemize}

	
%------------------------------------------------
\section{Forces appliquées à l’avion}\label{sec:section1}
	
	De manière générale, plusieurs forces s’appliquent à un avion : le poids, les forces aérodynamiques (portance et traînée), la poussée ainsi que la force de frottement avec le sol pendant la phase de roulage.
	
	\begin{enumerate}
		
		\item \textbf{Le poids ($W$)} correspond au poids total de l’avion, c’est-à-dire la somme du poids à vide, du poids des passagers et du carburant.
		
		\item \textbf{La portance ($L$)} est une force aérodynamique générée par la différence de pression entre l’intrados et l’extrados de l’aile. Elle permet à l’avion de décoller.
		
		\item \textbf{La traînée ($D$)} est une force aérodynamique qui s’oppose au mouvement de l’avion.
		
		\item \textbf{La poussée ($T$)} est la force produite par les moteurs permettant à l’avion d’accélérer.
		
		\item \textbf{La force de frottement ($F_r$)} est due au contact entre les roues de l’avion et la piste. Elle diminue progressivement à mesure que la portance augmente et devient nulle au moment du décollage.
		
	\end{enumerate}
	
	\vspace{1cm}
	
	
	
	%============================================================================================================
	%PARAMETRE INFLUENCANT LA DISTANCE DE DECOLLAGE 
	%================================================================================================================
	
	
	
	%------------------------------------------------
	\section{Paramètres influençant la distance de décollage}
	
		La distance de décollage dépend de plusieurs paramètres :
		
		\begin{itemize}
			\item la masse totale de l’avion ;
			\item la poussée des moteurs ;
			\item la densité de l’air ;
			\item la vitesse du vent ;
			\item la longueur et la nature de la piste ;
			\item les caractéristiques aérodynamiques de l’avion.
		\end{itemize}
		
		
		
		
		
		
	
	
	
	
	
	
	
	
	
	
	
	
	
	
	
	
	
	






